% ============================================================================
% Doc. 268 -- CMB acoustic peaks from the T4 lattice structure
%
% Purpose: Step-by-step derivation of the CMB peak ratios from the
%          T4 topology of FFGFT. Only what was actually computed and
%          checked. No claims without basis.
%
% Author: Johann Pascher
% Date: June 2026
% Status: Ratio match = retrodiction (Step 4/17.1). Forward content = geometric
%         anharmonicity (codimension-1 J0, Step 17). P29 partial, P30 open.
%         Naive Bessel projection insufficient.
% ============================================================================

\section*{Overview}

This document examines the relation between the T$^4$ lattice structure of
FFGFT and the observed CMB peak ratios. All numbers are directly
reproducible. What is a consistency check (backward) and what is a genuine
forward prediction is clearly separated from Step~17 on -- a key point that
was blurred in earlier versions of this document.

\textbf{Two separate statements that must not be conflated:}
\begin{itemize}
  \item \textbf{(A) Ratio finding (consistency check, backward).}
        The values $|n|^2 = 3, 18, 42, 78$ produce ratios
        $1\,:\,2{.}449\,:\,3{.}742\,:\,5{.}099$, which agree with the observed
        CMB peaks ($1\,:\,2{.}441\,:\,3{.}682\,:\,5{.}091$) to $<2\%$.
        \emph{However:} the sequence $\{1,6,14,26\}$ (before the factor~3) was
        \emph{read off} the measured CMB ratios (Step~4), not produced forward
        from T$^4$. This is a successful \emph{retrodiction}, not a
        parameter-free prediction (see Step~17.1).
  \item \textbf{(B) Anharmonicity finding (forward, new).}
        Forward, T$^4$ produces either a dense mode spectrum or -- for
        symmetric resonance -- the \emph{harmonic} series $1:2:3:4:5$. The CMB
        is \emph{anharmonic} ($1:2{.}44:3{.}68:5{.}09$). This anharmonicity
        follows \emph{forward and parameter-free} from the codimension-1
        projection in the fractal geometry ($D_f-1 \approx 2 \Rightarrow$
        Bessel index $\nu\approx0 \Rightarrow J_0$ membrane), matching the
        observed series to $\sim5\%$ -- with no mode selection (Step~17). This
        is the conceptually cleaner statement and the actual progress of this
        document.
\end{itemize}

\textbf{Honest summary:} the ratio match is real but found backward. The
forward content lies in the \emph{anharmonicity} as a geometric effect of the
flat codimension-1 projection. The rigorous selection mechanism for exactly
$\{3,18,42,78\}$ remains open (P29/P30, Doc.~190); in particular $|n|^2=30$ is
\emph{not} excluded by the current factor-3 filter (Step~17.1).

% ============================================================================
\section{Step 1: The T\textsuperscript{4} state space}
% ============================================================================

On the T$^4$ torus every mode has an integer winding number in
four directions:
\[
  \vec{n} = (n_1, n_2, n_3, n_4) \in \mathbb{Z}^4
\]
The radius (the quantum number) of a mode is:
\[
  |\vec{n}| = \sqrt{n_1^2 + n_2^2 + n_3^2 + n_4^2}
\]

The lowest levels, each with an example tuple:

\begin{center}
\begin{tabular}{rrlr}
\toprule
$|n|^2$ & $|n|$ & Example tuple & Degeneracy $g$ \\
\midrule
1 & 1{.}0000 & $(-1,0,0,0)$ & 8 \\
2 & 1{.}4142 & $(-1,-1,0,0)$ & 24 \\
3 & 1{.}7321 & $(-1,-1,-1,0)$ & 32 \\
4 & 2{.}0000 & $(-2,0,0,0)$ & 24 \\
5 & 2{.}2361 & $(-2,-1,0,0)$ & 48 \\
6 & 2{.}4495 & $(-2,-1,-1,0)$ & 96 \\
7 & 2{.}6458 & $(-2,-1,-1,-1)$ & 64 \\
8 & 2{.}8284 & $(-2,-2,0,0)$ & 24 \\
9 & 3{.}0000 & $(-3,0,0,0)$ & 104 \\
10 & 3{.}1623 & $(-3,-1,0,0)$ & 144 \\
\bottomrule
\end{tabular}
\end{center}

\textbf{Important:} The degeneracy $g(|n|^2)$ is \emph{not} monotonic.
At $|n|^2 = 6$ it jumps to 96 -- twelve times more than at $|n|^2 = 1$.
This follows from the Jacobi formula for the number of representations as
a sum of four squares: $g(n) = 8 \sum_{d | n,\, 4\nmid d} d$.

% ============================================================================
\section{Step 2: Thermal occupation (Bose-Einstein)}
% ============================================================================

The energy of a mode with radius $r = |n|$ is (from the calibration
$L_\mathrm{tor} = 1/(k_B T_\mathrm{CMB})$, Doc.~025):
\[
  E_r = r \cdot k_B T_\mathrm{CMB}
\]
The thermal occupation by Bose-Einstein with $\mu = 0$:
\[
  \langle N(r) \rangle = \frac{1}{e^r - 1}
\]
The spectral density -- how much energy in total sits in level $r$:
\[
  \boxed{I(r) = g(r) \cdot \langle N(r) \rangle \cdot r}
\]

Numerical values for the lowest levels:

\begin{center}
\begin{tabular}{rrrrrr}
\toprule
$|n|^2$ & $|n|$ & $g$ & $\langle N\rangle$ & $I = g\cdot N\cdot |n|$ & Local max.? \\
\midrule
1  & 1{.}0000 & 8   & 0{.}582 & 4{.}66  & \\
2  & 1{.}4142 & 24  & 0{.}321 & 10{.}90 & \\
3  & 1{.}7321 & 32  & 0{.}215 & 11{.}91 & \textbf{Yes} \\
4  & 2{.}0000 & 24  & 0{.}157 & 7{.}51  & \\
5  & 2{.}2361 & 48  & 0{.}120 & 12{.}84 & \\
6  & 2{.}4495 & 96  & 0{.}094 & 22{.}22 & \textbf{Yes} \\
7  & 2{.}6458 & 64  & 0{.}076 & 12{.}93 & \\
8  & 2{.}8284 & 24  & 0{.}063 & 4{.}26  & \\
9  & 3{.}0000 & 104 & 0{.}052 & 16{.}35 & \\
10 & 3{.}1623 & 144 & 0{.}044 & 20{.}13 & \textbf{Yes} \\
\bottomrule
\end{tabular}
\end{center}

The complete list of local maxima up to $|n|^2 = 80$:

\begin{center}
\begin{tabular}{rrrr}
\toprule
Peak & $|n|^2$ & $|n|$ & $I$ \\
\midrule
1 & 3  & 1{.}7321 & 11{.}91 \\
2 & 6  & 2{.}4495 & 22{.}22 \\
3 & 10 & 3{.}1623 & 20{.}13 \\
4 & 14 & 3{.}7417 & 17{.}45 \\
5 & 18 & 4{.}2426 & 19{.}30 \\
6 & 22 & 4{.}6904 & 12{.}52 \\
7 & 26 & 5{.}0990 & 10{.}52 \\
8 & 30 & 5{.}4772 & 13{.}25 \\
9 & 42 & 6{.}4807 & 7{.}64  \\
10 & 66 & 8{.}1240 & 2{.}77 \\
11 & 78 & 8{.}8318 & 1{.}73 \\
\bottomrule
\end{tabular}
\end{center}

% ============================================================================
\section{Step 3: Projection onto multipole moments}
% ============================================================================

The observable multipole moment $\ell$ is proportional to the
wavenumber of the mode:
\[
  \ell \propto |\vec{n}| = \sqrt{|n|^2}
\]
Thus the \emph{ratios} of the peak $\ell$ values are directly the
ratios of the $\sqrt{|n|^2}$ values -- independent of the absolute
scale (which contains $D_A$ and $L_\mathrm{tor}$):
\[
  \frac{\ell_k}{\ell_1} = \frac{\sqrt{|n|^2_k}}{\sqrt{|n|^2_1}}
\]

% ============================================================================
\section{Step 4: The decisive finding}
% ============================================================================

The CMB ratios \emph{squared} give directly the required
$|n|^2$ ratios:

\begin{center}
\begin{tabular}{rrrrl}
\toprule
CMB peak & $\ell_\mathrm{obs}$ & $\ell/\ell_1$ & $(\ell/\ell_1)^2$ & nearest $|n|^2$ \\
\midrule
1 & 220  & 1{.}000 & 1{.}00  & 1 \\
2 & 537  & 2{.}441 & 5{.}96  & \textbf{6} \\
3 & 810  & 3{.}682 & 13{.}56 & \textbf{14} \\
4 & 1120 & 5{.}091 & 25{.}92 & \textbf{26} \\
5 & 1440 & 6{.}545 & 42{.}84 & \textbf{43} \\
\bottomrule
\end{tabular}
\end{center}

If the \emph{first} CMB peak lies at $|n|^2 = X$, then the
further ones lie at $X \cdot 6$, $X \cdot 14$, $X \cdot 26$, \ldots

For $X = 3$: the peaks lie at $|n|^2 = 3, 18, 42, 78$.

\begin{tcolorbox}[colback=red!4!white, colframe=red!55!black,
  title={Honest reading of Step~4 (direction of inference)}]
The sequence $\{1,6,14,26\}$ here arises by \emph{squaring the measured} CMB
ratios and rounding to the nearest integer -- i.e.\ \emph{from the data}, not
forward from T$^4$. Step~5 then checks whether $3\cdot\{1,6,14,26\}$ are
genuine local maxima of the T$^4$ spectrum (they are). This is a
\textbf{consistency check (retrodiction)}, not a prediction. The inference
runs $\text{data}\to\{1,6,14,26\}\to{\times3}$, not $\text{T}^4\to\text{data}$.
The genuine forward content -- anharmonicity as a geometric effect -- is
developed in Step~17.
\end{tcolorbox}

% ============================================================================
\section{Step 5: Comparison with the T\textsuperscript{4} peaks}
% ============================================================================

Are $|n|^2 = 3, 18, 42, 78$ genuine local maxima in the T$^4$ spectrum?

\begin{center}
\begin{tabular}{rrrrl}
\toprule
$|n|^2$ & $g$ & $I$ & $I_\mathrm{left}$ & Local maximum? \\
\midrule
3  & 32   & 11{.}91 & 10{.}90 & \textbf{Yes} \\
18 & 312  & 19{.}30 & 9{.}77  & \textbf{Yes} \\
42 & 768  & 7{.}64  & 3{.}57  & \textbf{Yes} \\
78 & 1344 & 1{.}73  & 1{.}71  & \textbf{Yes} \\
\bottomrule
\end{tabular}
\end{center}

All four are genuine local maxima. The ratios:

\begin{center}
\begin{tabular}{rllll}
\toprule
 & Peak 1 & Peak 2 & Peak 3 & Peak 4 \\
\midrule
$|n|^2$ & 3 & 18 & 42 & 78 \\
$\sqrt{|n|^2}$ & 1{.}7321 & 4{.}2426 & 6{.}4807 & 8{.}8318 \\
Ratio to Peak 1 & 1{.}000 & 2{.}449 & 3{.}742 & 5{.}099 \\
CMB ratio & 1{.}000 & 2{.}441 & 3{.}682 & 5{.}091 \\
Deviation & 0{.}0\% & $+0{.}3\%$ & $+1{.}6\%$ & $+0{.}2\%$ \\
\bottomrule
\end{tabular}
\end{center}

\textbf{Deviation $< 2\%$ -- with no free parameter.}

The ratios follow solely from the Jacobi degeneracy structure
of the T$^4$ lattice and Bose-Einstein statistics.

% ============================================================================
\section{Step 6: What determines the absolute scale}
% ============================================================================

The ratios are scale-independent. The absolute position of the
first peak ($\ell_1 \approx 220$) depends on:
\[
  \ell_1 = \frac{\pi \cdot D_A}{L_\mathrm{tor} / \sqrt{3}}
\]
with the angular distance $D_A$ (external parameter, R20) and
$L_\mathrm{tor} = 1/(k_B T_\mathrm{CMB})$ (internal, from $\xi$).
This sets $A \approx 303{.}6$ in the trumpet formula (Doc.~267).

% ============================================================================
\section{Step 7: What is still open (R29)}
% ============================================================================

The T$^4$ spectrum has peaks at $|n|^2 = 3, 6, 10, 14, 18, 22, 26, 30, \ldots$
The CMB shows only the subset $3, 18, 42, 78$.

Between these lie strong peaks -- in particular $|n|^2 = 6$ with
$I = 22{.}22$, the largest value in the whole spectrum.

\textbf{Question (R29):} What selection rule explains why
$|n|^2 = 6, 10, 14, 22, 26, 30$ are not visible as CMB peaks?

\textbf{This question is answered in Steps~8 and~9.} The key is that
the observer measures in 3D space (plus time-unfolding): they see not
the full T$^4$ radius $|n|^2$ but only the spatial projection
$k_\mathrm{obs}^2 = (n_1^2 + n_2^2 + n_3^2)/L^2$. The most symmetric
modes give the factor~3, so that the visible series
$3\cdot\{1,6,14,26\} = \{3,18,42,78\}$ arises (Doc.~190, R29).

What \emph{rigorously} remains is only the full spherical projection
of the T$^4$ modes onto the angular power spectrum $C_\ell$ with
spherical Bessel functions (Doc.~190, R30) -- a feasible computation,
no fundamental obstacle. The ratios themselves are computed and stand.

% ============================================================================
\section{Summary}
% ============================================================================

\begin{center}
\begin{tabular}{p{4cm}p{8cm}}
\toprule
\textbf{Step} & \textbf{Result} \\
\midrule
1. T$^4$ state space &
  Winding numbers $\vec{n} \in \mathbb{Z}^4$,
  Jacobi degeneracy $g(|n|^2)$ non-monotonic \\
\addlinespace
2. Bose-Einstein &
  $I(r) = g(r) \cdot \frac{1}{e^r-1} \cdot r$,
  local maxima at $|n|^2 = 3, 6, 10, 14, 18, \ldots$ \\
\addlinespace
3. Projection &
  $\ell \propto \sqrt{|n|^2}$,
  ratios scale-independent \\
\addlinespace
4. CMB ratios &
  $(\ell/\ell_1)^2 \approx 1, 6, 14, 26$ for $|n|^2 = 1, 6, 14, 26$
  (or $\times 3$ for $|n|^2 = 3, 18, 42, 78$) \\
\addlinespace
5. Agreement &
  $\sqrt{3}:\sqrt{18}:\sqrt{42}:\sqrt{78} = 1:2{.}449:3{.}742:5{.}099$
  vs. CMB $1:2{.}441:3{.}682:5{.}091$, deviation $< 2\%$ \\
\addlinespace
6. Absolute scale &
  $\ell_1 \approx 220$ from $D_A$ (external) and $L_\mathrm{tor}$
  (internal, from $\xi$) \\
\addlinespace
7. R29/R30 &
  Selection rule for the subset $3, 18, 42, 78$: answered by the
  3D observer (Step~8/9); rigorous $C_\ell$ projection (R30)
  outstanding \\
\bottomrule
\end{tabular}
\end{center}

\noindent\textbf{Related:} Doc.~025 ($T_\mathrm{CMB}$ from $\xi$),
Doc.~050 (T$^4$ topology, spin structure),
Doc.~190 (R14a path-lengthening, R20 $R_H$ as external parameter),
Doc.~203 (recursion operator $\mathcal{R}$),
Doc.~267 (cosmological degeneracy, trumpet formula).

\noindent\emph{Python script:} \texttt{ffgft\_cmb\_t4\_peaks.py}

% ============================================================================
\section{Step 8: Where do factor 3 and $k^2$ come from?}
% ============================================================================

\subsection*{8.1 Where does the $k^2$ come from?}

The $^2$ in $k_\mathrm{eff}^2$ does not come from the factor~3 --
it comes from the wave equation itself.
The wave equation for a massless field ($m=0$, photon):
\[
  \left(\frac{\partial^2}{\partial t^2} - c^2\,\Delta\right)\phi = 0
\]
The Laplace operator $\Delta$ contains \emph{second} derivatives.
Inserting the Fourier ansatz $\phi \sim e^{i(\mathbf{k}\cdot\mathbf{x}-\omega t)}$
gives:
\[
  -\omega^2 + c^2 k^2 = 0
  \quad\Rightarrow\quad
  \omega^2 = c^2 k^2
  \quad\Rightarrow\quad
  E = \hbar\omega = \hbar c\,k
\]
The $^2$ is universal -- it holds for any wave equation,
independent of FFGFT.

\subsection*{8.2 Where does the factor 3 come from?}

The observer lives in \textbf{3 spatial dimensions + time-unfolding}.
They perceive the T$^4$ torus through their 3 spatial directions.
The fourth torus direction is the time-unfolding -- it is not
another spatial direction but the mechanism through which the
observer exists at all.

Concretely: a T$^4$ mode with winding number
$\vec{n} = (n_1, n_2, n_3, n_4)$ appears to the 3D observer
with the observable 3D wavenumber:
\[
  k_\mathrm{obs}^2 = \frac{n_1^2 + n_2^2 + n_3^2}{L_\mathrm{tor}^2}
\]
The $n_4$ component is ``hidden'' in the time-unfolding.

For the T$^4$ ground mode $(1,1,1,0)$ -- all three spatial directions
equal, $n_4 = 0$ -- one has:
\[
  k_\mathrm{obs}^2 = \frac{1^2+1^2+1^2}{L_\mathrm{tor}^2}
  = \frac{3}{L_\mathrm{tor}^2}
\]
This is exactly the \textbf{factor~3}: it counts the three
observable spatial dimensions.

The fractal spatial dimension $D_f = 3 - \xi \approx 3$
(Doc.~051) is consistent with this: $D_f$ is the number of
spatial dimensions, slightly modified by $\xi$. The correction
$\xi/3 \approx 4{.}4 \times 10^{-5}$ is negligible.

\subsection*{8.3 The observable series}

The observable multipole moment is
$\ell \propto k_\mathrm{obs} \cdot D_A$:
\[
  \ell^2 \propto k_\mathrm{obs}^2 \cdot D_A^2
  = \frac{n_1^2+n_2^2+n_3^2}{L_\mathrm{tor}^2} \cdot D_A^2
\]
For the dominant T$^4$ ground peaks at $|n|^2 = 1, 6, 14, 26$
-- all with the property that $n_1^2+n_2^2+n_3^2 = 3\cdot|n_\mathrm{min}|^2$
for the most symmetric mode -- one has:
\[
  \ell^2 \propto 3 \cdot |n|^2
\]
The visible peaks $\{3, 18, 42, 78\} = 3\cdot\{1, 6, 14, 26\}$:

\begin{center}
\begin{tabular}{rllll}
\toprule
 & Peak 1 & Peak 2 & Peak 3 & Peak 4 \\
\midrule
T$^4$ ground peak $|n|^2$ & 1 & 6 & 14 & 26 \\
$\times 3$ (spatial dim.) & 3 & 18 & 42 & 78 \\
$\ell \propto \sqrt{3\cdot|n|^2}$ & 1{.}732 & 4{.}243 & 6{.}481 & 8{.}832 \\
Ratio to Peak~1 & 1{.}000 & 2{.}449 & 3{.}742 & 5{.}099 \\
CMB ratio & 1{.}000 & 2{.}441 & 3{.}682 & 5{.}091 \\
Deviation & $0{.}0\%$ & $+0{.}3\%$ & $+1{.}6\%$ & $+0{.}2\%$ \\
\bottomrule
\end{tabular}
\end{center}

\subsection*{8.4 What is still open}

The statement ``$n_1^2+n_2^2+n_3^2 = 3\cdot|n|^2$ for the dominant
modes'' holds exactly only for modes of the form $(n,n,n,0)$.
The general distribution over all modes of an $|n|^2$ level
gives on average $\langle n_1^2+n_2^2+n_3^2\rangle = \tfrac{3}{4}|n|^2$
(isotropic T$^4$ distribution), not $3\cdot|n|^2$.

The exact selection mechanism -- why the 3D observer prefers
precisely the coincidence peaks
$3\cdot\{1,6,14,26\}=\{3,18,42,78\}$ -- is not yet rigorously derived.
\textbf{The physical motivation} (3 spatial dimensions,
observer in 3D + time-unfolding) is clear.
The full spherical projection of the T$^4$ modes onto $C_\ell$
with spherical Bessel functions would show it --
an open computation, no fundamental obstacle (R30).

% ============================================================================
\section{Step 9: Complete derivation -- overview}
% ============================================================================

The CMB peak ratios follow from three steps:

\begin{center}
\begin{tabular}{c p{2.1cm} p{6.3cm} p{3.8cm}}
\toprule
Step & Input & Mechanism & Output \\
\midrule
1 & $\xi = 1/7500$ &
  Jacobi degeneracy $g(|n|^2)$, Bose-Einstein $N(r)$ &
  T$^4$ peaks at $|n|^2 = 1, 6, 14, 26$ \\
\addlinespace
2 & $D = 3$ spatial dim. &
  Observer in 3D + time-unfolding:
  $k_\mathrm{obs}^2 = (n_1^2+n_2^2+n_3^2)/L^2$ &
  Factor~3 in $\ell^2 \propto 3\cdot|n|^2$ \\
\addlinespace
3 & Steps 1+2 &
  $\ell \propto \sqrt{3\cdot|n|^2}$ &
  Ratios $1:2{.}449:3{.}742:5{.}099$ \\
\bottomrule
\end{tabular}
\end{center}

Observed: $1:2{.}441:3{.}682:5{.}091$. Deviation $< 2\%$.

\textbf{The $^2$} in $k_\mathrm{obs}^2$ comes from the wave equation
(second derivatives of the Laplace operator) -- universal, not an FFGFT specialty.

\textbf{The factor~3} comes from the 3 observable spatial dimensions --
the observer in 3D + time-unfolding sees only $n_1^2+n_2^2+n_3^2$,
not the full T$^4$ radius $|n|^2$.

\textbf{Open step:} The exact selection mechanism --
why the observer prefers the most symmetric modes
$(n_1^2+n_2^2+n_3^2 \approx 3\cdot|n_\mathrm{min}|^2)$ --
is not yet rigorously derived.

\textbf{Important reading (see Step~17):} The sequence $\{1,6,14,26\}$ listed
as the ``Output'' of Step~1 is \emph{not} produced forward from T$^4$; it is
read off the measured CMB ratios (Step~4). This ``overview'' therefore
describes the \emph{retrodiction}. The genuine forward content -- the
anharmonicity from the codimension-1 projection -- is in Step~17.

% ============================================================================
\section{Summary (updated)}
% ============================================================================

\begin{center}
\begin{tabular}{p{3.5cm}p{8.5cm}}
\toprule
\textbf{Step} & \textbf{Result} \\
\midrule
1. T$^4$ state space &
  Jacobi degeneracy $g(|n|^2)$ non-monotonic \\
\addlinespace
2. Bose-Einstein &
  Local maxima at $|n|^2 = 1, 6, 14, 26, \ldots$ \\
\addlinespace
3. Observer in 3D &
  $k_\mathrm{obs}^2 = (n_1^2+n_2^2+n_3^2)/L^2$,
  factor~3 from 3 spatial dimensions;
  $^2$ from wave equation (Laplace) \\
\addlinespace
4. CMB series &
  $3 \cdot \{1, 6, 14, 26\} = \{3, 18, 42, 78\}$,
  all genuine T$^4$ peaks \\
\addlinespace
5. Ratios &
  $1:2{.}449:3{.}742:5{.}099$ vs.\ CMB $1:2{.}441:3{.}682:5{.}091$,
  deviation $< 2\%$ \\
\addlinespace
6. Absolute scale &
  $D_A$ external (R20), $L_\mathrm{tor}$ internal from $\xi$ \\
\addlinespace
7. Still open &
  Rigorous selection mechanism:
  why $n_1^2+n_2^2+n_3^2 = 3\cdot|n_\mathrm{min}|^2$
  for the dominant modes? \\
\bottomrule
\end{tabular}
\end{center}

\noindent\textbf{Related:} Doc.~025 ($T_\mathrm{CMB}$ from $\xi$),
Doc.~050/051 (T$^4$ topology, Dirac structure),
Doc.~190 (R14a, R20, R29, R30),
Doc.~203 (recursion operator $\mathcal{R}$),
Doc.~267 (cosmological degeneracy).

% ============================================================================
\section{Step 10: The physical intuition -- double resonator}
% ============================================================================

\subsection*{10.1 Open and closed pipe}

In classical acoustics the boundary conditions determine which modes
resonate:

\begin{center}
\begin{tabular}{p{3.8cm}p{3.4cm}p{5.2cm}}
\toprule
\textbf{System} & \textbf{Boundary condition} & \textbf{Mode ratios} \\
\midrule
Closed pipe & both ends fixed & $1:2:3:4:5$ (all) \\
Half-open pipe & one end open & $1:3:5:7$ (odd only) \\
Circular membrane (2D) & rim fixed & $1:2{.}30:3{.}60:4{.}90$ \\
\midrule
\textbf{Fractal T$^4$} & \textbf{two coupled systems} &
  $\mathbf{1:2{.}449:3{.}742:5{.}099}$ \\
\bottomrule
\end{tabular}
\end{center}

The pure T$^4$ (integer dimension $D=4$, fully compact)
corresponds to a \emph{closed} system -- all modes resonate.
The fractal T$^4$ (dimension $D_f = 3-\xi < 4$) has a
qualitatively different resonance structure.

\subsection*{10.2 The fractal T$^4$ as a double resonator}

The analogy to the half-open pipe (which produces a factor~2) is
tempting -- but it is not complete. The fractal T$^4$ is
neither a simply open nor a simply closed system. It acts
as a \textbf{double resonator}: two superimposed resonance structures
whose coincidences determine the visible peaks.

\begin{center}
\begin{tabular}{p{3cm}p{5cm}p{5cm}}
\toprule
 & \textbf{System 1} & \textbf{System 2} \\
 & T$^4$ lattice & 3D observer projection \\
\midrule
Resonances & $|n|^2 = 1, 6, 14, 26, \ldots$ &
  Projected onto $n_1^2+n_2^2+n_3^2$ \\
Condition & Local maximum in $I(r)$ &
  $k_\mathrm{obs}^2 = (n_1^2+n_2^2+n_3^2)/L^2$ \\
Origin & Jacobi degeneracy $g(|n|^2)$ &
  3 spatial dimensions ($D_f = 3-\xi \approx 3$) \\
\bottomrule
\end{tabular}
\end{center}

\textbf{Only the coincidence resonances are visible} -- modes where
both systems resonate simultaneously:
\[
  |n|^2_\mathrm{visible}
  = 3 \cdot |n|^2_\mathrm{lattice}
  = 3 \cdot \{1,\,6,\,14,\,26\}
  = \{3,\,18,\,42,\,78\}
\]

\subsection*{10.3 Why not all coincidences are visible}

The double resonance condition yields more candidates than just
$\{3, 18, 42, 78\}$:

\begin{center}
\begin{tabular}{rrrrl}
\toprule
$|n|^2$ & $|n|^2/3$ & Peak in Sys.~1? & $I$ & CMB peak? \\
\midrule
3  & 1  & no & 11{.}91 & \textbf{Yes} (strongest in range) \\
18 & 6  & \textbf{Yes} & 19{.}30 & \textbf{Yes} \\
30 & 10 & \textbf{Yes} & 13{.}25 & No (weaker) \\
42 & 14 & \textbf{Yes} & 7{.}64  & \textbf{Yes} \\
54 & 18 & \textbf{Yes} & 4{.}54  & No (weaker) \\
66 & 22 & \textbf{Yes} & 2{.}77  & No (weaker) \\
78 & 26 & \textbf{Yes} & 1{.}73  & \textbf{Yes} \\
\bottomrule
\end{tabular}
\end{center}

The peaks $|n|^2 = 30, 54, 66$ satisfy the double resonance condition
but are thermally weaker ($I$ smaller) than $|n|^2 = 18, 42, 78$.
In each range the \emph{strongest} coincidence peak prevails:

\begin{itemize}
  \item Range $|n|^2 \in [1, 20]$:\quad
        Strongest coincidence peak: $|n|^2 = 18$ ($I = 19{.}3$) \\
        \hphantom{Range $|n|^2 \in [1, 20]$:\quad}
        $|n|^2 = 30$ would have $I = 13{.}2$ -- suppressed.
  \item Range $|n|^2 \in [20, 50]$:\quad
        Strongest coincidence peak: $|n|^2 = 42$ ($I = 7{.}6$) \\
        \hphantom{Range $|n|^2 \in [20, 50]$:\quad}
        $|n|^2 = 30$ is stronger ($I = 13{.}2$), but not a CMB peak --
        it lies between the coincidences.
  \item Range $|n|^2 \in [50, 90]$:\quad
        Strongest coincidence peak: $|n|^2 = 78$ ($I = 1{.}7$).
\end{itemize}

\subsection*{10.4 The physical meaning}

The transition from $D = 4$ (pure T$^4$, integer compact) to
$D_f = 3 - \xi$ (fractal T$^4$) corresponds to the transition from a
fully closed to a \emph{fractally partly-open} system:

\begin{itemize}
  \item \textbf{$D = 4$ (purely compact):} All four torus directions
        oscillate on equal footing. Every $|n|^2$ can resonate.
        This corresponds to a resonator with $4$ equivalent,
        closed directions.

  \item \textbf{$D_f = 3 - \xi$ (fractal):} The fourth direction is
        no longer integer-compact but fractally granulated.
        It behaves like a direction with a \emph{continuous
        spectrum} -- analogous to the open side of a half-open pipe.
        Decisive is: the observer in 3D space sees only the three
        spatial components $n_1^2+n_2^2+n_3^2$ (the fourth is
        time-unfolding). This produces the factor~3 in the
        observable projection and selects the
        coincidence resonances.
\end{itemize}

\textbf{In brief:} The fractal T$^4$ is neither a closed nor an
open resonator -- it is a \emph{double resonator} with coupled
systems. The CMB peaks mark the coincidences where
T$^4$ lattice geometry and 3D observer projection simultaneously
satisfy resonance conditions.

\subsection*{10.5 Comparison of resonator types}

\begin{center}
\begin{tabular}{p{3.5cm}p{3.5cm}p{5cm}}
\toprule
\textbf{Analogy} & \textbf{Mechanism} & \textbf{Selects} \\
\midrule
Closed pipe & one boundary condition & all integers \\
Half-open pipe & one boundary condition & odd only \\
Circular membrane & 2D geometry & Bessel zeros \\
\midrule
\textbf{Fractal T$^4$} &
  \textbf{coincidence of two} \newline \textbf{resonator systems} &
  $D_f \cdot \{$lattice peaks$\}$ \\
\bottomrule
\end{tabular}
\end{center}

The CMB peak series $1:2{.}449:3{.}742:5{.}099$ is thus the
resonance signature of a fractally partly-open, four-dimensional
resonance space -- not of a plasma, not of a simple membrane,
but of the specific combination of T$^4$ topology and
fractal dimension $D_f = 3 - \xi$.

\subsection*{10.6 What follows from the peaks: two filters}

The discrete, sharp peaks follow from two properties of the same
complete T$^4$ acting together -- seen from two reference points:

\begin{center}
\begin{tabular}{p{2.2cm}p{4.6cm}p{4.6cm}}
\toprule
 & \textbf{Filter~1: compactness} & \textbf{Filter~2: fractality} \\
 & (T$^4$ as a whole, from outside) & (from inside, decompactified) \\
\midrule
Effect & produces \emph{discrete} spectrum &
  weights the modes unequally \\
\addlinespace
forces & \emph{existence} of sharp peaks &
  \emph{location} of the dominant peaks \\
\addlinespace
Origin & finite, closed topology &
  Jacobi degeneracy $g(|n|^2)$ + Bose-Einstein \\
\addlinespace
Result & uniform mode lattice &
  peaks at $|n|^2 = 3, 6, 10, 14, \ldots$ \\
\bottomrule
\end{tabular}
\end{center}

\textbf{The two are not two objects} but two reference points on
the same complete T$^4$: compact-finite from outside (Filter~1),
fractally granulated inside down to $L_0 = \xi\cdot\ell_P$ (Filter~2).
The fractal non-closure is therefore not a \emph{second} property
contradicting compactness, but the inside view of the same closed
torus -- a matter of reference point (inside/outside, R20).

\textbf{Only both filters together} give the observed structure:
\begin{itemize}
  \item Filter~1 alone: uniform discrete lattice -- every
        $|n|^2$ equal, no preferred peaks.
  \item Filter~2 alone: needs a discrete lattice to act on --
        no lattice without compactness.
  \item Filter~1 $+$ Filter~2: dominant peaks at specific $|n|^2$,
        which after 3D projection (Step~8) give the CMB series.
\end{itemize}

\textbf{What can be deduced:} The discrete peaks argue \emph{for}
a compact (finite) universe and \emph{against} a purely-infinite-open
one -- a purely open space would give a continuum without sharp
peaks. This is the strongest structural statement that can be drawn
from the peaks.

\textbf{Honest caveat:} This conclusion is FFGFT-internal. The peaks
do \emph{not} force spatial compactness -- $\Lambda$CDM produces the
same discrete peaks from a \emph{temporal} bound (the sound horizon
at recombination) in a spatially infinite universe. The peaks alone
do not distinguish spatial compactness (FFGFT) from a temporal bound
($\Lambda$CDM); this is the cosmological degeneracy (Doc.~267).

% ============================================================================
\subsection*{10.7 What the peak ratios state model-free}

This section considers the measured peak positions purely with the
methods of spectral / frequency analysis -- without any T$^4$
assumption. The question: what can be read from the bare numbers
$\ell \approx 220, 537, 810, 1120, 1450$, and where do typical
misinterpretations from differencing and sampling limits lurk?

\subsubsection*{10.7.0 Clarification: ``acoustic'' is misleading}

The CMB peaks are historically called ``acoustic peaks''. The term
sends anyone who hears it for the first time in the wrong direction,
because there is \emph{nothing} to do with sound:

\begin{itemize}
  \item \textbf{No sound:} there is no audible signal, no pressure
        wave in air.
  \item \textbf{No oscillation in time:} the CMB is a frozen still
        image, not a time signal. Nothing is measured in hertz.
  \item \textbf{No direct wave measurement:} the telescope measures
        the temperature as a function of \emph{sky direction} -- a map
        of hot and cold spots. What is evaluated is the \emph{spatial
        distribution on the sky}, not an oscillation.
\end{itemize}

The word ``acoustic'' comes from the $\Lambda$CDM \emph{explanatory
model}: there the spots are interpreted as frozen pressure waves in
the photon-baryon plasma \emph{before} recombination. That is a
model-dependent interpretation of the \emph{mechanism} -- it already
sits in the name, before anything is measured. The \emph{measured
quantity} itself is neutral: a spatial interference pattern on the
sky, described by the angular power spectrum $C_\ell$.

\textbf{Two different ``frequencies'' -- do not confuse them.} At the
CMB two completely different things carry the name ``frequency'':

\begin{center}
\begin{tabular}{p{1cm}p{5.3cm}p{5.3cm}}
\toprule
 & \textbf{(A) Photon frequency} & \textbf{(B) Spot pattern} \\
\midrule
What & colour/energy of the light & angular structure of the spots \\
Value & $\sim 160$\,GHz & $\ell = 220, 537, \ldots$ \\
Unit & hertz (temporal) & dimensionless / angle (spatial) \\
measured as & radiometer (photon energy) & angle on the sky \\
\bottomrule
\end{tabular}
\end{center}

The peak ratios $1:2{.}44:3{.}68:\ldots$ belong to \textbf{(B)},
the \emph{spatial} angular pattern. They have nothing to do with (A),
the photon frequency. Converting the multipoles into a hertz frequency
would be artificial and misleading -- it would translate a spatial
pattern into a fictitious temporal oscillation that does not exist.

\textbf{Convention in this document:} we avoid ``acoustic'' where
possible and write \emph{angular-spectrum peaks} or simply
\emph{CMB peaks}. The resonator analogy (open/closed pipe, double
resonator) remains a purely \emph{mathematical} analogy for a discrete
eigenmode spectrum -- not a claim that there was sound. Whether the
resonances arise from pressure waves ($\Lambda$CDM) or from geometric
lattice modes (FFGFT) is exactly the open question that the name
``acoustic'' prematurely decides in favour of $\Lambda$CDM.

\subsubsection*{10.7.1 First and second difference: the curvature is real}

The spacings of successive peaks are \emph{not} constant:
\[
  \Delta\ell = 317,\; 273,\; 310,\; 330.
\]
A pure harmonic series (integer overtones $1:2:3:4:5$)
would have constant spacings. The ratios $1:2{.}44:3{.}68:5{.}09:6{.}59$
grow \emph{more slowly} than $1:2:3:4:5$ -- the characteristic curve is
\textbf{concave}. This curvature is much larger than the measurement
uncertainty of the peak positions ($\sim 1$--$3\,\%$ for broad peaks).
It is therefore \emph{not} a differencing artefact but a real feature
of the spectrum.

\textbf{Conclusion 1:} integer overtones (a simple harmonic
fundamental) are excluded.

\subsubsection*{10.7.2 The trap: ``missing fundamental'' when reading linearly}

From the practice of frequency analysis it is known: if one measures
only overtones and the fundamental is missing, the spacings suggest a
fundamental that is not actually excited. Reading the CMB peaks
linearly ($\ell_n = A + B\,n$), the fit gives:
\[
  A = -72, \qquad B = 297.
\]
The offset $A$ is \textbf{negative}. The linear continuation thus
demands a ``fundamental'' at $\ell < 0$ -- physically meaningless.
This is exactly the missing-fundamental artefact: whoever reads the
peaks as ``overtones of a shifted fundamental'' constructs a
fundamental mode that does not exist.

\textbf{Conclusion 2:} the linear reading (integer overtones plus
offset) is a differencing artefact and is ruled out.

\subsubsection*{10.7.3 The fundamental in the lattice sense is reconstructible}

If one instead checks whether the peaks are
$\sqrt{\text{integer}}$ multiples of a common fundamental, i.e.
$\ell = c\cdot\sqrt{|n|^2}$ with $|n|^2 = 3, 18, 42, 78$, then
$c = \ell/\sqrt{|n|^2}$ is nearly constant:
\[
  c = 127{.}0;\; 126{.}6;\; 125{.}0;\; 126{.}8
  \quad(\text{scatter } 1{.}6\,\%).
\]
There thus exists a common \textbf{lattice fundamental}
$c \approx 126$. This is the most parsimonious reading of the
curvature: \emph{one} parameter (the scale $c$) suffices, since the
curvature itself sits parameter-free in the fixed lattice numbers
$3, 18, 42, 78$.

\begin{center}
\begin{tabular}{p{4.5cm}p{3.5cm}p{3.5cm}}
\toprule
 & \textbf{harmonic} & \textbf{lattice ($\sqrt{|n|^2}$)} \\
\midrule
Fundamental definition & $\ell_n = f_0\cdot n$ & $\ell = c\cdot\sqrt{|n|^2}$ \\
Constancy test & $\ell/n = 220\ldots290$ & $\ell/\sqrt{|n|^2} \approx 126$ \\
Parameters & 2 (offset $+$ slope) & 1 (scale $c$ only) \\
Offset & $A = -72$ (unphysical) & none needed \\
Result & fails & $c \approx 126$ \\
\bottomrule
\end{tabular}
\end{center}

\textbf{Conclusion 3:} the fundamental is reconstructible -- but only
in the lattice sense ($\ell = c\sqrt{|n|^2}$), not the harmonic one.
The value $c \approx 126$ corresponds to the mode $|n|^2 = 1$.

\subsubsection*{10.7.4 The reconstructed fundamental is not the first peak}

The lattice fundamental $c \approx 126$ corresponds to $\ell \approx 126$
for $|n|^2 = 1$. The first \emph{visible} peak, however, lies at
$|n|^2 = 3$, i.e. $\ell \approx 219$. There is thus a \textbf{real but
peak-invisible} fundamental mode:
\[
  |n|^2 = 1 \;\to\; \ell \approx 126
  \quad(\text{fundamental, not a peak}),
\]
\[
  |n|^2 = 3 \;\to\; \ell \approx 219
  \quad(\text{first visible peak}).
\]
Remarkably, $\ell \approx 126$ lies \emph{in the middle of the
observed range}, but there is no peak there, rather the rising
Sachs-Wolfe plateau. This is consistent with the observer mechanism
(Step~8): the mode $|n|^2=1$ is the single-axis mode $(1,0,0,0)$, which
does \emph{not} satisfy the symmetric 3D projection condition; only
$(1,1,1,0)$ with $|n|^2=3$ projects isotropically and becomes the
first peak.

\textbf{Conclusion 4:} there is a real, lower-lying fundamental mode at
$\ell \approx 126$ that acts as a scale-setter but, owing to missing
symmetry, does not appear as a peak. This is a weakly falsifiable
consistency condition: the fundamental $c$ reconstructed from the upper
peaks must match the place where the spectrum passes from the plateau
into the first peak.

\subsubsection*{10.7.5 No size of the universe from the oscillations}

No \emph{absolute size} can be read from the angular oscillations --
neither the size of the structures nor that of the universe. The
measured angle of a peak is
\[
  \theta = \frac{L_\mathrm{structure}}{D_A},
\]
i.e. only the \emph{ratio} of physical structure size
$L_\mathrm{structure}$ and distance $D_A$ (angular distance to the
source surface). The same angle fits ``large structure, far away''
just as well as ``small structure, nearby'':

\begin{center}
\begin{tabular}{p{4cm}p{4cm}}
\toprule
\textbf{assumed $D_A$} & \textbf{$\to L_\mathrm{structure}$ (at $\theta\approx0{.}8°$)} \\
\midrule
$14\,000$\,Mpc & $\approx 200$\,Mpc \\
$100$\,Mpc & $\approx 1{.}4$\,Mpc \\
$1$\,Mpc & $\approx 0{.}01$\,Mpc \\
\bottomrule
\end{tabular}
\end{center}

To turn the angle into an absolute length one needs $D_A$ -- and
$D_A$ is not a measured quantity but a \textbf{model output}:
\begin{itemize}
  \item $\Lambda$CDM: $D_A = \int c/H(z)\,dz$ -- requires $H_0$,
        $\Omega_m$, $\Omega_\Lambda$, the recombination $z_*$.
  \item FFGFT: $D_A = R_H\ln(1+z_*)$ -- requires $R_H$, the external
        parameter (R20), which does not follow from $\xi$ alone.
\end{itemize}

In both cases the distance comes from the model, not from the
measurement. The angular measurement is \textbf{scaleless}: it shows
the \emph{form} of the pattern (the ratios), but not its absolute
size -- like a photograph without a scale bar shows the form of an
object, not its size. This is exactly the same limit as for the peak
ratios: the \emph{form} follows from the geometry, the \emph{absolute
scale} needs an external parameter (R20).

\textbf{Why the resonator method does not help here.} For a real
resonator in front of us -- an organ pipe, say -- the size does sit in
the fundamental: measuring the fundamental wavelength $\lambda_0$
\emph{in metres} and knowing the medium (sound speed
$c_\mathrm{medium}$) gives the length directly, $L = \lambda_0/2$. This
method needs two things, both of which are absent at the CMB:

\begin{itemize}
  \item \textbf{A wavelength in metres} -- not just an angle. The organ
        pipe is in front of us; we measure its length directly. The
        CMB ``resonator'' is unreachably far away; we see only its
        angle on the sky, not its length. The step from angle to length
        needs $D_A$ -- the unknown model output.
  \item \textbf{A known medium.} For the organ pipe $c_\mathrm{medium}$
        is a known constant. In the $\Lambda$CDM plasma picture the
        ``sound speed'' $c_s(z) = c/\sqrt{3(1+R(z))}$ is not a constant
        but a \emph{function computed from the model}: it requires the
        baryon and photon densities as functions of time, the expansion
        history and the recombination state. Here the ``medium'' is not
        a substance one has in front of one and whose $c$ one measures
        -- it arises only \emph{indirectly} from the assumed curvature
        and expansion. ``Denser in the past'' is already a statement
        about the geometry, not about an independently existing medium.
\end{itemize}

Thus in the $\Lambda$CDM picture the back-calculation is triply
model-dependent: the ``wavelength'' is only an angle, the ``medium''
$c_s(z)$ is computed from density, expansion and metric, and the
distance $D_A$ is likewise a model output. None of the three
ingredients is a model-free measurement; the ``medium'' is even the
most derived quantity. FFGFT in turn needs \emph{no} medium and
\emph{no} sound speed -- the peaks are eigenmodes of a geometry, not
sound in a substance -- but does need $R_H$ as an external parameter
for the absolute scale (R20). Not assumption-free, but a different
\emph{kind} of assumption: a geometric scale instead of a whole plasma
and expansion history.

\subsubsection*{10.7.6 What is certain and what is not -- the sampling limit}

\textbf{Certainly readable (model-free):}
\begin{itemize}
  \item No integer overtones -- the curvature is real, not an artefact.
  \item The curvature is concave (ratios grow sub-linearly).
  \item A discrete mode lattice with $\sqrt{|n|^2}$ scaling is the
        most parsimonious reading (one parameter, $c \approx 126$).
  \item There is no harmonic fundamental ``below'' the first peak;
        the linear fit there is a missing-fundamental artefact.
  \item No \emph{absolute size} (of the universe or the structures)
        can be read -- the angle gives only $L/D_A$, and $D_A$ is a
        model output.
\end{itemize}

\textbf{Not readable (limit of the method):}
\begin{itemize}
  \item Whether the concave curvature stems from a
        $\sqrt{\text{integer lattice}}$ (FFGFT) or from acoustic
        $\Lambda$CDM dispersion -- both produce the same form
        (cosmological degeneracy, Doc.~267).
  \item Fine differences between curved laws: the CMB peaks are
        \emph{broad} (half-width $\Delta\ell \sim 100$--$200$), the
        positions determinable only to $\sim 1$--$3\,\%$. This suffices
        to \emph{exclude} integer overtones, not to distinguish the
        concave laws.
\end{itemize}

\textbf{Conclusion of this section:} pure frequency analysis -- without
any T$^4$ assumption -- yields three robust statements: (i)~no harmonic
overtone series, (ii)~the curvature is real and concave, (iii)~a
lattice fundamental $c \approx 126$ exists and corresponds to a real
but invisible $|n|^2{=}1$ mode. These statements are obtained
independently of the FFGFT derivation and support it from the
measurement side, without removing the cosmological degeneracy (FFGFT
vs.\ $\Lambda$CDM).

% ============================================================================
\section{Step 11: No fit -- what is derived and what remains open}
% ============================================================================

The agreement of the peak ratios to $< 2\%$ does \emph{not} come from
a fit. This section shows explicitly what comes from where.

\subsection*{11.1 What follows from $\xi$ alone -- and what does not}

\begin{tcolorbox}[colback=red!4!white, colframe=red!55!black,
  title={Correction relative to earlier versions}]
Earlier versions claimed the sequence $\{1,6,14,26\}$ followed ``from $\xi$
alone as local maxima''. This is \emph{not} correct. The genuine local maxima
of the full thermally weighted T$^4$ spectrum lie at
$|n|^2 = 3, 6, 10, 14, 18, 22, 26, 30, \ldots$ (dense). Of these, $6,14,26$
are local maxima, but $1$ is \emph{not}, and $10$ (a local maximum) is missing
from the sequence. So $\{1,6,14,26\}$ is \emph{not} a natural forward subset --
it was read off the data (Step~4). What follows forward from T$^4$ is cleanly
separated in Step~17.
\end{tcolorbox}

\begin{center}
\begin{tabular}{p{4cm}p{8cm}}
\toprule
\textbf{Quantity} & \textbf{Origin} \\
\midrule
T$^4$ spectrum maxima at \newline $|n|^2 = 3,6,10,14,18,22,26,30,\ldots$ &
  Local maxima of $I(r) = g(r)\cdot N(r)\cdot r$;
  $g$ from the Jacobi formula (exact), $N$ from Bose-Einstein.
  Input: only $\xi = 1/7500$. \emph{Forward, but dense.} \\
\addlinespace
Harmonic backbone $1:2:3:4:5$ &
  Symmetric resonance $(n,n,n)$: $k^2_\mathrm{obs}=3n^2$.
  \emph{Forward}, but harmonic -- ``too clean''. \\
\addlinespace
Anharmonicity (stretch) &
  Codimension-1 projection $D_f-1\approx2 \Rightarrow J_0$ Bessel.
  \emph{Forward, parameter-free} (Step~17). \\
\addlinespace
Selection of exactly $\{1,6,14,26\}$ &
  \emph{Read off the CMB data} (Step~4), not forward. \\
\bottomrule
\end{tabular}
\end{center}

\textbf{Forward:} the dense spectrum, the harmonic backbone and the geometric
anharmonicity. \textbf{Backward:} the concrete selection of the four visible
peaks.

\subsection*{11.2 What is not yet rigorously derived}

\begin{center}
\begin{tabular}{p{4cm}p{4cm}p{4cm}}
\toprule
\textbf{Assumption} & \textbf{Status} & \textbf{Consequence} \\
\midrule
$L_\mathrm{tor} = 1/(k_BT_\mathrm{CMB})$ &
  Physically motivated, not rigorously proved &
  Affects only the absolute scale, \emph{not} the ratios \\
\addlinespace
$k_\mathrm{obs}^2 = (n_1^2+n_2^2+n_3^2)/L^2$ &
  Observer in 3D + time-unfolding,\newline $n_4$ = time-unfolding &
  Affects the factor~3; physically motivated, mathematically open \\
\addlinespace
$D_A$ (angular distance) &
  External parameter $R_H$ needed (R20) &
  Affects only $\ell_1 \approx 220$, not the ratios \\
\bottomrule
\end{tabular}
\end{center}

\subsection*{11.3 What the ratios prove and what they do not}

The ratios $1:2{.}449:3{.}742:5{.}099$ agree with the CMB to
$< 2\%$ \emph{without} free parameters. This shows:

\begin{itemize}
  \item The T$^4$ geometry and the 3D observer projection
        \emph{automatically} produce a peak structure of the right
        order of magnitude.
  \item The deviation of $< 2\%$ lies within the measurement
        uncertainty of the raw CMB peak positions
        (without a $\Lambda$CDM template fit).
\end{itemize}

It does \emph{not} prove:

\begin{itemize}
  \item That the selection mechanism (why the 3D observer prefers
        precisely the most symmetric modes with
        $n_1^2+n_2^2+n_3^2 = 3\cdot|n_\mathrm{min}|^2$) follows
        rigorously from the spherical $C_\ell$ projection
        (still to be shown, R30).
  \item That no other theory can produce the same ratios.
  \item That FFGFT correctly describes the entire CMB
        (amplitudes, polarisation, early-phase physics open).
\end{itemize}

\textbf{The status is:} an unfitted result from $\xi$ that
matches the observed ratios to $< 2\%$.
This is a genuine test -- falsifiable by more precise model-free
measurement of the raw peak positions.

% ============================================================================
\section{Step 12: Schwarzschild connection and $R_H$}
% ============================================================================

\subsection*{12.1 The two Schwarzschild radii in FFGFT}

FFGFT has two scales that are simultaneously Schwarzschild radii
(Doc.~190, R4; Doc.~182):

\begin{center}
\begin{tabular}{p{2.5cm}p{4cm}p{4cm}l}
\toprule
\textbf{Scale} & \textbf{Definition} & \textbf{Schwarzschild mass} & \textbf{Check} \\
\midrule
$L_0 = \xi\cdot\ell_P$ &
  sub-Planck fundamental length &
  $M_0 = \xi\cdot m_P/2$ &
  $r_S(M_0) = L_0$ \\
\addlinespace
$R_H = \hbar c/E_H$ &
  Hubble horizon &
  $M_U = R_H c^2/(2G)$ &
  $r_S(M_U) = R_H$ \\
\bottomrule
\end{tabular}
\end{center}

Both checks numerically:
\[
  r_S(M_0) = \frac{2G\cdot\xi\cdot m_P/2}{c^2}
  = \xi\cdot\frac{G\,m_P}{c^2}
  = \xi\cdot\ell_P = L_0
  \quad\checkmark
\]
\[
  r_S(M_U) = \frac{2G\cdot M_U}{c^2}
  = \frac{2G}{c^2}\cdot\frac{R_H c^2}{2G}
  = R_H
  \quad\checkmark
\]

\subsection*{12.2 The conversion: $R_H$ from the two Schwarzschild conditions}

Dividing the two Schwarzschild formulas:
\[
  \frac{R_H}{L_0}
  = \frac{2G\,M_U/c^2}{2G\,M_0/c^2}
  = \frac{M_U}{M_0}
\]

Thus:
\[
  \boxed{R_H = L_0 \cdot \frac{M_U}{M_0}
  = \xi\cdot\ell_P\cdot\frac{M_U}{\xi\cdot m_P/2}
  = \frac{2\,\ell_P\cdot M_U}{m_P}}
\]

This is a pure conversion -- no model, no assumption.
It shows: $R_H$ follows directly from $L_0$ and the mass ratio
$M_U/M_0$.

\subsection*{12.3 What this means}

$L_0 = \xi\cdot\ell_P$ is fully derived from $\xi$.
$M_0 = \xi\cdot m_P/2$ likewise.

$R_H$ is thus fully determined \emph{once} $M_U$ is known:
\[
  R_H = \frac{2\,\ell_P}{m_P}\cdot M_U
\]

In Doc.~182 $M_U$ is obtained via $E_H = E_0\cdot\xi^{41/4}$.
The Schwarzschild conversion shows the same $R_H$ from another
side -- as a direct consequence of the mass scale of the universe.

\subsection*{12.4 Open question: $M_U$ independently from $\xi$?}

The connection $R_H = 2\ell_P M_U/m_P$ is exact.
Open is: can $M_U$ be derived from the T$^4$ mode structure
\emph{without} presupposing $R_H$?

If so, then $R_H$ would be a genuine prediction from $\xi$ alone --
without needing the exponent $41/4$ as input.
This is an open derivation, no fundamental obstacle.

% ============================================================================
\section{Step 13: The measurement problem -- what FFGFT can do, what $\Lambda$CDM cannot}
% ============================================================================

\subsection*{13.1 The exponent 41/4 is equivalent to $H_0$}

From $E_H = E_0 \cdot \xi^{41/4}$ and $H_0 = E_H/\hbar$ it follows:
\[
  \frac{41}{4} = \frac{\ln(E_0/(\hbar H_0))}{\ln(1/\xi)}
\]
The exponent $41/4$ is thus \textbf{equivalent to measuring $H_0$}.
Numerically:

\begin{center}
\begin{tabular}{lll}
\toprule
$H_0$ [km/s/Mpc] & Source & Exponent \\
\midrule
$66{.}19$ & FFGFT (from $41/4$) & $10{.}2492$ \\
$67{.}40$ & Planck/CMB ($\Lambda$CDM) & $10{.}2472$ \\
$73{.}00$ & SH0ES (distance ladder) & $10{.}2382$ \\
\bottomrule
\end{tabular}
\end{center}

The difference between the exponents is smaller than $0{.}02$ --
all $H_0$ measurements give the same exponent to $\sim 0{.}1\%$.

\subsection*{13.2 The measurement problem}

There is no model-free measurement of $H_0$.
Every measurement presupposes model assumptions:

\begin{center}
\begin{tabular}{p{3.5cm}p{5.5cm}p{3cm}}
\toprule
\textbf{Method} & \textbf{Model assumption} & \textbf{Result} \\
\midrule
CMB peaks ($\Lambda$CDM) &
  expansion, FLRW metric, recombination &
  $67{.}4$ km/s/Mpc \\
\addlinespace
Distance ladder (SH0ES) &
  calibrated standard candles, local geometry &
  $73{.}0$ km/s/Mpc \\
\addlinespace
FFGFT ($\xi^{41/4}$) &
  static universe, T$^4$ topology,\newline
  exponent $41/4$ from theory &
  $66{.}2$ km/s/Mpc \\
\bottomrule
\end{tabular}
\end{center}

In FFGFT $H_0$ means something different than in $\Lambda$CDM:
not the expansion rate, but the reciprocal of the
fractal path-lengthening scale $R_H = c/H_0$.
What $\Lambda$CDM measures as expansion, FFGFT interprets
as fractal path-lengthening $z(d)$ (Doc.~190, R14a).

\subsection*{13.3 What FFGFT can do, what $\Lambda$CDM cannot}

\begin{center}
\begin{tabular}{p{5cm}p{4.5cm}p{4.5cm}}
\toprule
\textbf{Quantity} & \textbf{$\Lambda$CDM} & \textbf{FFGFT} \\
\midrule
$H_0$ &
  Free parameter\newline (determined from data) &
  Derived from $\xi^{41/4}$\newline (if $41/4$ derived) \\
\addlinespace
CMB peak ratios &
  From $r_s(z^*)$ and $D_A(z^*)$\newline (needs $H_0$, $\Omega_b$, \ldots) &
  From T$^4$ geometry\newline (no free parameter) \\
\addlinespace
Absolute peak position $\ell_1$ &
  From $H_0$, $z^*$, $r_s$ &
  From $R_H$ (from $\xi^{41/4}$)\newline + open connection R20 \\
\bottomrule
\end{tabular}
\end{center}

\textbf{The structural difference:}
$\Lambda$CDM needs $H_0$ as input -- measured, not derived.
FFGFT writes $H_0 = 66{.}2$\,km/s/Mpc as $\xi^{41/4}$.
If $41/4$ follows from the T$^4$ geometry (still open),
$H_0$ is a \emph{prediction} -- not a free parameter.

\textbf{The remaining open problem:}
The exponent $41/4$ is not derived from the T$^4$ geometry in any
document. It is the only place in FFGFT cosmology where an external
input is needed -- equivalent to $H_0$.
As long as this derivation is missing, FFGFT and $\Lambda$CDM stand
on equal footing in this respect: both need a cosmological scale
as input.

% ============================================================================
\section{Step 14: The universal Compton-Schwarzschild relation}
% ============================================================================

\subsection*{14.1 Every length has two associated masses}

To every length $R$ belong two masses:
\[
  M_c(R) = \frac{\hbar}{R\,c}
  \quad\text{(}R\text{ is its Compton wavelength)}
\]
\[
  M_s(R) = \frac{R\,c^2}{2G}
  \quad\text{(}R\text{ is its Schwarzschild radius)}
\]

Their product is \emph{independent of $R$}:
\[
  \boxed{M_c(R)\cdot M_s(R)
  = \frac{\hbar}{Rc}\cdot\frac{Rc^2}{2G}
  = \frac{\hbar c}{2G}
  = \frac{m_P^2}{2}}
\]

\subsection*{14.2 Application to the electron}

Equivalently in lengths: to every mass $m$ belong the Compton
wavelength $\lambda = \hbar/(mc)$ and the Schwarzschild radius
$r_S = 2Gm/c^2$, with
\[
  \lambda \cdot r_S = \frac{\hbar}{mc}\cdot\frac{2Gm}{c^2}
  = \frac{2G\hbar}{c^3} = 2\,\ell_P^2
\]
Thus:
\[
  \sqrt{\lambda_e \cdot r_S(e)} = \sqrt{2}\,\ell_P
\]
The two scales of the electron lie \emph{geometrically symmetric}
around the Planck length -- independent of the electron mass.

\begin{center}
\begin{tabular}{lll}
\toprule
\textbf{Scale} & \textbf{Value} & \textbf{Meaning} \\
\midrule
$r_S(e) = 2Gm_e/c^2$ & $1{.}35\times10^{-57}$\,m & Schwarzschild (bottom) \\
$\ell_P$ & $1{.}62\times10^{-35}$\,m & Planck (middle) \\
$\lambda_e = \hbar/(m_e c)$ & $3{.}86\times10^{-13}$\,m & Compton (top) \\
\bottomrule
\end{tabular}
\end{center}

\subsection*{14.3 Application to $R_H$}

The same relation holds for the Hubble radius:
\[
  M_c(R_H) = \frac{\hbar}{R_H c} = \frac{E_H}{c^2}
  \quad\text{(tiny)}
\]
\[
  M_s(R_H) = \frac{R_H c^2}{2G} = M_U
  \quad\text{(mass of the universe)}
\]
with
\[
  M_c(R_H)\cdot M_s(R_H) = \frac{E_H}{c^2}\cdot M_U
  = \frac{m_P^2}{2}
  \quad\checkmark
\]

\subsection*{14.4 What this relation achieves and what it does not}

\textbf{It achieves:} an elegant internal consistency.
The Schwarzschild mass of the universe $M_U$ and the
cosmological energy $E_H$ are connected by $M_U \cdot E_H/c^2 = m_P^2/2$
-- exact, without free parameter. The two Schwarzschild scales
$L_0$ (bottom) and $R_H$ (top) are not arbitrary bounds
but Schwarzschild radii of real masses ($M_0 = \xi m_P/2$ and $M_U$).

\textbf{It does not achieve:} a determination of $R_H$.
The relation $M_c \cdot M_s = m_P^2/2$ is a universal identity --
it holds for \emph{every} length and therefore gives no new information
about the numerical value of $R_H$. Both masses $M_c(R_H)$ and $M_s(R_H)$
are defined via $R_H$; their product is tautologically $m_P^2/2$.

To fix $R_H$ one needs an \emph{independent} input:
the cosmological energy $E_H = E_0\cdot\xi^{41/4}$ (Doc.~182).
The electron fixes the Planck scale via
$\sqrt{\lambda_e r_S(e)} = \sqrt{2}\ell_P$, but not $R_H$.

\textbf{The remaining open problem stays:}
The exponent $41/4$ -- equivalent to the cosmological energy scale $E_H$ --
is the only input that FFGFT needs for the absolute cosmological scale.
The Compton-Schwarzschild relation explains the \emph{structure}
of the scale hierarchy, but not the \emph{value} of $E_H$.

% ============================================================================
\section{Step 15: The T$^4$ length is a Schwarzschild radius}
% ============================================================================

\subsection*{15.1 The fundamental identity (Doc.~010, 066, 067)}

The FFGFT documents define the T0 length of an object as:
\[
  r_0 = 2Gm \quad\text{(natural units, } c=1\text{)}
\]
Doc.~067 explicitly states: the factor~2 results from the
relativistic field structure and \emph{coincides exactly with the
Schwarzschild radius}. Doc.~010 writes $r_0 = 2GE$
(energy instead of mass), Doc.~066 as the ratio $\xi = r_0/\ell_P = 2Gm/\ell_P$.

\textbf{This means:} every characteristic T0 length \emph{is}
a Schwarzschild radius -- not analogous, but identical.

\subsection*{15.2 Why different formulas in the documents?}

It is the same formula in different unit systems:

\begin{center}
\begin{tabular}{p{3cm}p{3.5cm}p{4.5cm}}
\toprule
\textbf{Document} & \textbf{Formula} & \textbf{Unit system} \\
\midrule
SI & $r_S = 2Gm/c^2$ & metres, kg, seconds \\
Doc.~015, 029 & $r_S = 2M$ & natural ($G=c=1$) \\
Doc.~010 & $r_0 = 2GE$ & $c=1$, energy instead of mass \\
Doc.~066 & $r_0 = \xi\cdot\ell_P$ & Planck ($\ell_P=1$) \\
\bottomrule
\end{tabular}
\end{center}

For the minimal mass $M_0 = (\xi/2)\,m_P$ \emph{all} give the same value:
\[
  \frac{2G M_0}{c^2}
  = \frac{2G}{c^2}\cdot\frac{\xi}{2}\,m_P
  = \xi\cdot\frac{G\,m_P}{c^2}
  = \xi\cdot\ell_P = L_0 = 2{.}155\times10^{-39}\,\text{m}
\]
The apparent variety is only the SI conversion, which shows or hides
factors $c^2$ and $G$. As soon as one computes in SI units, $c^2$ and
$G$ appear explicitly; in natural units they vanish. The physics is
identical in all cases.

\subsection*{15.3 ``a Schwarzschild radius'' vs.\ ``the Schwarzschild radius''}

In standard physics one speaks of \emph{the} Schwarzschild radius
of an object -- the one length at which this mass becomes a black hole.

In FFGFT the T0 length is \emph{always} a Schwarzschild radius --
but \emph{which one} depends on the mass:

\begin{center}
\begin{tabular}{lll}
\toprule
\textbf{Mass} & \textbf{Schwarzschild radius} & \textbf{Role} \\
\midrule
Electron $m_e$ & $1{.}35\times10^{-57}$\,m & particle scale \\
$M_0 = (\xi/2)m_P$ & $L_0 = 2{.}16\times10^{-39}$\,m & fundamental length \\
Proton & $2{.}48\times10^{-54}$\,m & particle scale \\
$M_U$ (universe) & $R_H = 1{.}41\times10^{26}$\,m & Hubble horizon \\
\bottomrule
\end{tabular}
\end{center}

\textbf{Therefore:} $L_0$ \emph{is a} Schwarzschild radius
(namely that of $M_0$) -- but \emph{not the} Schwarzschild radius
as such, because there is no single universal one.
Likewise $R_H$ is \emph{a} Schwarzschild radius (that of $M_U$),
not \emph{the} one.

This is the scale relativity from Doc.~182: every system carries
its own torus with its own characteristic scale
$R_\mathrm{torus}(m) = \hbar/(mc)$. The universal lower bound $L_0$
holds for all; a universal upper bound does not exist --
$R_H$ is only the upper bound of the cosmological sector.

\subsection*{15.4 What this means for $R_H$}

$R_H$ \emph{coincides exactly} with the Schwarzschild radius of $M_U$ --
this is correctly stated in the documents and is not an approximation.
But: because $r_0 = 2Gm$ is the \emph{definition} of the T0 length,
the agreement is necessary, not surprising. The T0 length
of a mass and its Schwarzschild radius are the same object.

The identity does not fix $R_H$ independently: one still needs
$M_U$ (resp.\ $E_H = E_0\cdot\xi^{41/4}$). What the identity achieves
is the \emph{physical interpretation}: $R_H$ is not an abstract
horizon but the Schwarzschild radius of the total mass of the
static universe -- the universe sits at its own
Schwarzschild bound.

% ============================================================================
\section{Step 16: Why the Schwarzschild relation does not derive $R_H$}
% ============================================================================

\subsection*{16.1 The tempting but false conclusion}

From two true formulas a fallacy arose:
\begin{align}
  \text{(A)} \quad & \xi = L_0/\ell_P, \quad L_0 = \xi\cdot\ell_P \\
  \text{(B)} \quad & r_S(m)\cdot r_C(m) = 2\,\ell_P^2 \quad\text{(universal)}
\end{align}

Both are correct. The error arose through a tacit assumption: that in
relation (B) one may replace the electron mass $m_e$ by the universe
mass $M_U$ \emph{while keeping $\lambda_e$}, thus obtaining
$R_H = r_S(M_U)$.

\subsection*{16.2 The exact error}

Relation (B) connects two quantities of the \emph{same} mass. For the
electron:
\[
  \underbrace{r_S(m_e)}_{1{.}35\times10^{-57}\,\text{m}}
  \cdot
  \underbrace{\lambda_e}_{3{.}86\times10^{-13}\,\text{m}}
  = 2\,\ell_P^2
\]
Both factors belong to $m_e$.

Replacing $m_e \to M_U$ replaces \emph{both} factors:
\[
  \underbrace{r_S(M_U)}_{R_H \,=\, 1{.}41\times10^{26}\,\text{m}}
  \cdot
  \underbrace{r_C(M_U)}_{3{.}7\times10^{-96}\,\text{m}}
  = 2\,\ell_P^2
\]
The Compton wavelength of $M_U$ is $3{.}7\times10^{-96}$\,m --
\textbf{not} $\lambda_e$. The conflation was: $R_H$ (Schwarzschild
of the \emph{universe}) was put together with $\lambda_e$ (Compton of
the \emph{electron}) into a relation that holds only for a
single mass. In fact:
\[
  r_S(M_U)\cdot\lambda_e
  = 5{.}4\times10^{13}\,\text{m}^2
  \neq 2\,\ell_P^2 = 5{.}2\times10^{-70}\,\text{m}^2
\]
-- off by a factor of $10^{83}$.

\subsection*{16.3 Two separate systems}

Electron and universe are two different systems, each consistent in
itself:
\[
  \text{Electron:}\quad m_e \to \{r_S(e),\,\lambda_e\},
  \quad r_S(e)\cdot\lambda_e = 2\ell_P^2
\]
\[
  \text{Universe:}\quad M_U \to \{R_H,\,r_C(M_U)\},
  \quad R_H\cdot r_C(M_U) = 2\ell_P^2
\]
Between the systems there is \emph{no} bridge from this relation.

\subsection*{16.4 The Schwarzschild formula generates no information}

The deeper reason: $r_S = 2GM/c^2$ is a \emph{conversion}
between mass and length. As for any object:

\begin{center}
\begin{tabular}{p{4cm}p{8cm}}
\toprule
\textbf{Known} & \textbf{Follows} \\
\midrule
Mass $M$ & Schwarzschild radius $r_S = 2GM/c^2$ \\
Length $r_S$ & Mass $M = r_S c^2/(2G)$ \\
\textbf{neither} & \textbf{nothing} \\
\bottomrule
\end{tabular}
\end{center}

This holds for a black hole, the Sun, the electron and the
universe alike. To know $R_H$, one of the four equivalent
quantities must be independently known:
\[
  M_U \;\longleftrightarrow\; H_0 \;\longleftrightarrow\;
  E_H \;\longleftrightarrow\; R_H
\]
They form a ring of conversions -- but no entry point.

\subsection*{16.5 Where FFGFT supplies the input}

FFGFT supplies exactly one cosmological input:
\[
  E_H = E_0\cdot\xi^{41/4}, \qquad E_0 = \sqrt{m_e\,m_\mu}
\]
From this follows $R_H = \hbar c/E_H$ and everything else. The exponent
$41/4$ is equivalent to specifying $H_0$ and is not derived from the
T$^4$ geometry in any document (open problem).

\textbf{Conclusion:} The Schwarzschild identity ($L_0$, $R_H$ are
Schwarzschild radii) is a genuine structural property and supplies
the \emph{physical interpretation}. But it cannot \emph{generate}
the cosmological scale -- for that the one input $E_H$ is needed.
The earlier impression that the two formulas (A) and (B) would derive
$R_H$ rested on the tacit substitution $m_e \to M_U$, which conflates
two different masses.

% ============================================================================
\section{Step 17: Forward test -- anharmonicity from the codimension-1 projection}
% ============================================================================

This step cleanly separates what follows \emph{forward} (without looking at
the data) from T$^4$ and what is a \emph{backward} consistency check. It
documents an explicit forward investigation and is the conceptual advance over
Steps 1--11. All numbers here are reproducible with the accompanying scripts
(\texttt{python/Dok267\_268\_Skripte/}).

\subsection*{17.1 What the naive forward computation yields -- and what it does not}

Three forward computations, all without any look at the data:

\begin{enumerate}
  \item \textbf{Symmetric T$^4$, spherical $C_\ell$ Bessel projection}
        ($C_\ell = \sum_n N_\mathrm{BE}\,|j_\ell(k_\mathrm{obs}D_A)|^2$):
        does \emph{not} reproduce $\{3,18,42,78\}$. The result is a
        quasi-continuum dominated by the lowest modes (strongest peak at
        $\ell\approx122$). The naive Bessel projection is \emph{insufficient}
        (contrary to the earlier P30 wording ``a feasible computation, no
        fundamental obstacle'').
  \item \textbf{3+time (T$^3$ spatial spectrum):} local maxima at
        $|n|^2=2,5,9,14,17,21,26,\ldots$ -- dense, and \emph{not}
        $\{1,6,14,26\}$.
  \item \textbf{Symmetric spatial resonance $(n,n,n)$:} $k^2_\mathrm{obs}=3n^2$
        $\Rightarrow$ series $\{3,12,27,48,75\}$, ratios
        $\sqrt{3}:\sqrt{12}:\sqrt{27}:\ldots = 1:2:3:4:5$ (\emph{harmonic}).
\end{enumerate}

\textbf{Conclusion:} forward, T$^4$ produces either a dense spectrum or the
harmonic backbone $1:2:3:4:5$. The sparse, anharmonic series $\{3,18,42,78\}$
is \emph{not} a forward output; it arises only from the data-read selection
$\{1,6,14,26\}$.

\textbf{Concretely open ($|n|^2=30$, sharpened by O.~Teker):} $30=3\cdot10$,
and $|n|^2=10$ is a genuine local maximum ($I=20.1$). So $30$ satisfies the
factor-3 filter. Thermally it is even \emph{stronger} than two of the visible
peaks: $I(30)=13.2 > I(42)=7.6 > I(78)=1.7$. Yet the CMB shows no peak at
$|n|^2=30$ ($\ell\approx696$, in the trough between peaks 2 and 3). The
``strongest in range'' argument (earlier version) fails here. The selection
mechanism is thus \emph{not} supplied by the factor~3 alone (P29 only partial,
P30 open).

\subsection*{17.2 The actual forward finding: anharmonicity is geometric}

The harmonic backbone $1:2:3:4:5$ is ``too clean'' -- like the ideal overtones
of a stiffness-free string. Real strings are \emph{inharmonic} due to their
stiffness (stretched piano tuning). By analogy we \emph{expect} anharmonicity,
because spacetime is curved from our viewpoint and the projection is nonlinear
(fractal path lengthening).

Acoustic resonators show that anharmonicity is purely \emph{geometric}
(forward, no fit, no medium):

\begin{center}
\begin{tabular}{p{5.5cm}p{4.5cm}p{2cm}}
\toprule
\textbf{Resonator (dimension $D$)} & \textbf{Ratios} & \textbf{Dev.\ to CMB} \\
\midrule
Tube/sphere $j_0$ ($D=3$, $\nu=\tfrac12$) & $1:2:3:4$ (harmonic) & $-18$ to $-24\%$ \\
\textbf{Circular membrane $J_0$ ($D=2$, $\nu=0$)} & $\mathbf{1:2.30:3.60:4.90}$ & $\mathbf{-2}$ to $\mathbf{-6\%}$ \\
S$^3$ Laplacian & $1:1.63:2.24$ & $-33$ to $-48\%$ \\
S$^2$ (Y$_{\ell m}$) & $1:1.73:2.45$ & too dense \\
\midrule
\textbf{CMB observed} & $1:2.44:3.68:5.09$ & --- \\
\bottomrule
\end{tabular}
\end{center}

\textbf{Key fact:} the radial Bessel index is $\nu = D/2-1$. Anharmonicity is
\emph{maximal at $D=2$} ($\nu=0$, $J_0$) and vanishes at $D=1$ and $D=3$ (both
harmonic). The CMB series sits exactly at $D\approx2$ -- the dimension of
maximal anharmonicity.

\textbf{Important:} the match is to the \emph{flat} circular membrane ($J_0$),
\emph{not} to the curved 2-sphere (Y$_{\ell m}$, which gives $1:1.73:\ldots$,
wrong). The correct expectation is a \emph{flat bounded disk}, not a sphere --
because the CMB peaks come from the \emph{radial} projection, not from the
eigenvalues of the observation surface.

\subsection*{17.3 Derivation: flat disk from the fractal codimension-1}

The last-scattering surface is a \textbf{codimension-1 hypersurface} in the
fractal space:
\[
  D_\mathrm{eff} = D_f - 1 = (3-\xi) - 1 = 2 - \xi \approx 2.
\]
Radial modes in $D$ dimensions are zeros of $J_{D/2-1}$. For $D = 2-\xi$:
\[
  \nu = \frac{D}{2}-1 = \frac{2-\xi}{2}-1 = -\frac{\xi}{2}
      \approx -6.7\times10^{-5} \approx 0
  \;\Rightarrow\; J_0\text{ membrane}.
\]
This gives, forward, \textbf{parameter-free and without mode selection}, the
anharmonic series $1:2.30:3.60:4.90$ -- a match to $\sim5\%$. The anharmonicity
follows from the codimension-1 projection in the fractal geometry, exactly as
expected from the curvature.

\subsection*{17.4 The small residual -- honestly open}

The best fit lies at $\nu\approx-0.07$, i.e.\ effective dimension
$D\approx1.86$ -- just \emph{below} 2. The obvious physical source would be the
\emph{spectral dimension} $d_s$ (governing wave propagation on a fractal),
which is generally below the Hausdorff dimension: $d_s = 2D_f/d_w$ with walk
dimension $d_w>2$.

The explicit computation rules this out, however:
\begin{itemize}
  \item The FFGFT fractal is \emph{too weakly fractal}: with $D_f=3-\xi$,
        $d_w = 2+O(\xi)$, so $d_s(\text{surface})\approx2.00$. $\xi$ is
        $\sim1000\times$ too small to push $D$ down to $1.86$.
  \item For $d_s=1.86$ one would need $d_w\approx2.15$ (strong fractality),
        incompatible with $D_f=2-\xi$.
\end{itemize}
\textbf{The residual is not derivable from $\xi$.} Moreover it is of the order
of the measurement uncertainty: the pure $J_0$ prediction ($D=2$) is consistent
with peaks 3 and 4 within $\sim2\sigma$; only peak 2 sits at $\sim4\sigma$. The
``best fit $1.86$'' is thus largely measurement noise; pointing to an exact
value (e.g.\ $13/7$) would be numerology unsupported by the data quality.

\subsection*{17.5 Summary -- layered, honest structure}

\begin{center}
\begin{tabular}{p{5.5cm}p{6.5cm}}
\toprule
\textbf{Layer} & \textbf{Status} \\
\midrule
$D=3\to2$ (codimension-1, $J_0$) &
  \textbf{forward, parameter-free, derived}; $\sim5\%$ \\
Anharmonicity geometric &
  derived (maximal at $D=2$) \\
Ratio match $\{3,18,42,78\}$ &
  retrodiction (from data), \emph{not} prediction \\
$D=2\to1.86$ (sub-2) &
  not from $\xi$; largely measurement noise; open \\
Exclusion of $|n|^2=30$ &
  open (P29 partial, P30) \\
Peak-2 residual ($\sim4\sigma$) &
  small genuine open point, likely other physics \\
\bottomrule
\end{tabular}
\end{center}

\textbf{Defensible statement (instead of ``fit to $D=1.86$''):} the fractal
codimension-1 projection produces, forward and parameter-free, the $J_0$
membrane spectrum ($D\approx2$) and matches the CMB anharmonicity to $\sim5\%$
-- within the measurement errors. The harmonic backbone comes from the
symmetric T$^4$ resonance; the stretch from the geometry of observation. This
formulation is conceptually cleaner than the earlier mode selection because it
requires \emph{no} selection from the data, thereby removing the circularity
objection (O.~Teker). The rigorous selection mechanism for the discrete
individual peaks remains open.
